\section{The VEB Environment}
\label{sec:environment}

VEB models a B2B enterprise-sales engagement as a partially observable,
multi-agent sequential decision problem. A single \emph{seller} agent (the system
under test) interacts over a bounded calendar with a \emph{buying committee} of
LLM-simulated personas, each holding private information that the seller must earn
through skilled discovery. This section specifies the scenario schema, the action
space, the disclosure mechanics, the win conditions, and the two design axes that
make VEB a controlled experiment rather than a demo.

\subsection{Scenario schema}
\label{sec:schema}
Each scenario is a declarative specification (a versioned YAML file) with the
following components:

\begin{itemize}[leftmargin=1.4em,itemsep=2pt,topsep=2pt]
  \item \textbf{Seller brief and account.} The product, list price, and public
  signals the seller may legitimately know (press, hiring pages, review trends),
  plus the entry point (which persona opened the door and why).
  \item \textbf{Personas.} A buying committee, each with a \texttt{committee\_role}
  (champion candidate, economic buyer, technical gatekeeper, procurement), a
  \texttt{personality}, \texttt{pressures}, and \texttt{private\_notes} visible
  only to the simulator. Exactly one persona is the \texttt{economic\_buyer} who
  can commit budget.
  \item \textbf{Org chart and budget cycle.} The approval topology and the
  \emph{compelling event}---an operational or fiscal deadline that makes inaction
  costly (a peak season, a capital-planning lock, an auto-renewing incumbent
  contract).
  \item \textbf{Competitor.} An incumbent or status-quo alternative with explicit
  \texttt{strengths} and \texttt{weaknesses}, so competitive handling can be
  graded against ground truth rather than vibes.
  \item \textbf{Gated facts.} The private facts that constitute the deal's
  substance, each with a \texttt{gate} specifying the conversational move that
  unlocks it (Section~\ref{sec:gates}).
  \item \textbf{Events.} Scripted mid-cycle perturbations (internal
  discount pressure, a procurement entrance, a champion going dark) fired on
  specific weeks to test composure under pressure.
  \item \textbf{Win conditions.} The machine-checkable bar for a clean win
  (Section~\ref{sec:wincond}).
\end{itemize}

\subsection{Scenario construction and calibration}
\label{sec:construction}
Scenarios enter the library through a construction pipeline designed so that a
scenario is admitted only if it is \emph{provably winnable but not trivially so},
and only if its deal logic is internally coherent. This is what lets us treat the
declarative win predicate (Section~\ref{sec:wincond}) as ground truth rather than
an aspiration. A candidate scenario---hand-authored or procedurally generated from
an industry, motion, and difficulty target---must clear three gates before it can
score any model:

\begin{enumerate}[leftmargin=1.4em,itemsep=2pt,topsep=2pt]
  \item \textbf{Schema and internal consistency (deterministic).} Beyond schema
  validity we run structural checks that catch the ways a generated deal can be
  quietly broken: every gated fact's holder must be a real persona; at least one
  pain must be earnable through a \texttt{quantifying\_question} gate; scripted
  perturbations must fire inside the calendar; and difficulty-appropriate event
  pressure must actually be present (a difficulty-$3{+}$ deal must stack multiple
  events including internal discount pressure). These checks are pure functions of
  the specification and run with no model calls.
  \item \textbf{Reachability.} Because access to facts and to the economic buyer is
  gated behind unlock chains, a naively generated scenario can be
  \emph{unwinnable}---a required fact or the budget holder sitting behind a gate no
  legal move opens. We statically trace the unlock graph and reject any scenario in
  which a \texttt{required\_fact} holder or the required economic buyer is
  unreachable, so ``the deal is impossible'' can never masquerade as ``the model
  failed.''
  \item \textbf{Solvability calibration (reference sellers).} Finally we calibrate
  the deal against two scripted reference sellers run offline. A deliberately weak
  \emph{baseline} seller must \emph{not} win (if it does, the scenario is too easy
  and is rejected), while a \emph{disciplined} reference seller that executes the
  methodology correctly must win (if it cannot, the scenario is likely unwinnable).
  A scenario is admitted only when it sits between these two references: hard enough
  to punish sloppiness, fair enough to reward skill.
\end{enumerate}

An optional LLM self-critique pass flags softer coherence issues (a persona whose
stated pressures contradict the org chart, say) for human review, but it never
gates admission on its own---the load-bearing gates are all deterministic or
reference-seller-checked. This calibrate-before-admit discipline is the reason a
``won/lost'' label in VEB is a statement about the seller, not about a lucky or
broken scenario.

\subsection{Action space and calendar}
\label{sec:actions}
The seller acts over a multi-week calendar (typically six weeks) with a fixed
number of interaction slots per week and a per-call turn budget. Within this
budget the seller may place calls, send emails, request meetings, and---crucially
---\emph{plan privately} between touches, producing internal artifacts (a value
model, a stakeholder map, a mutual action plan) that are graded even when never
sent. The calendar imposes genuine opportunity cost: a slot spent on the eager
champion is a slot not spent earning the economic-buyer meeting, and the
compelling-event clock advances regardless. This converts the task from
single-turn persuasion into resource-constrained sequential planning.

\subsection{Gated disclosure}
\label{sec:gates}
The defining mechanic of VEB is that \emph{facts are earned, not asked for}. Each
gated fact carries a \texttt{gate} type---for example \texttt{quantifying\_question}
(the seller must drive to a quantified impact, not merely acknowledge a pain),
\texttt{process\_question} (elicit the real decision process and its owners),
\texttt{champion\_test} (make an ask with internal cost, such as an EB
introduction), \texttt{competitor\_probe} (map alternatives before a landmine is
revealed), and \texttt{eb\_meeting\_held} (a commitment available only inside a
real economic-buyer meeting). A gate that unlocks the economic buyer's persona
(\texttt{unlocks\_personas}) makes stakeholder access itself a gated achievement:
the seller cannot talk to the budget holder until it has earned the introduction
in the correct way. Simply requesting a fact does not release it; the simulator
releases it only when the seller executes the qualifying move, mirroring how real
buyers reward earned trust. This is what lets VEB measure discovery \emph{skill}
rather than discovery \emph{effort}.

\subsection{Win conditions}
\label{sec:wincond}
Each scenario declares a machine-checkable win bar with several conjuncts: a set
of \texttt{required\_facts} that must have been elicited (in the buyer's own
words), a minimum trust threshold (\texttt{min\_trust}), whether a real
economic-buyer meeting was held (\texttt{requires\_eb\_meeting}), whether a mutual
action plan with customer-confirmed dates exists (\texttt{requires\_map}), and a
maximum discount the seller may concede while still counting as a disciplined
close (\texttt{max\_discount\_pct}). The last conjunct is essential: it encodes
that \emph{closing by gifting margin is not winning}. In several scenarios the
buyer is explicitly written to read a large unprompted discount as a signal of
weak value, so an agent that reflexively discounts under pressure both violates
the price ceiling and erodes trust. Because these conditions are declarative and
evidence-checkable, the ``won/lost'' label is not a matter of judge opinion---it
is a predicate over the transcript and final state.

\subsection{Design axis I: difficulty tiers}
\label{sec:tiers}
Scenarios are stratified into three difficulty tiers, calibrated by the number of
gates, the hostility and size of the committee, the presence of veto-holding
gatekeepers, and the tightness of the price ceiling.

\begin{itemize}[leftmargin=1.4em,itemsep=2pt,topsep=2pt]
  \item \textbf{Tier 1 (easy, $n{=}\numTierEasy$).} Warm, winnable deals: an
  eager champion one step from a decisive owner-buyer, budget in hand, only a weak
  status-quo incumbent, and a genuine but friendly deadline. The test is
  \emph{disciplined closing}---earning the buyer meeting and holding price rather
  than fumbling a gimme or buying the yes with a discount. Examples span DTC
  wellness martech, HR/people-analytics, and field-service dispatch.
  \item \textbf{Tier 2 (mid, $n{=}\numTierMid$).} Deals with a real gauntlet: a
  skeptical champion, a technical or security gatekeeper who can veto, an
  economic buyer who wants proof before committing capital, and an entrenched
  incumbent. Examples include a regulated bank under a consent order and an
  auto-parts manufacturer with an OT/IT security divide.
  \item \textbf{Tier 3 (hard, $n{=}\numTierHard$).} Adversarial, high-stakes
  deals where urgency is a trap. Example: a national retailer six weeks after a
  credential-stuffing breach, where the motivated CISO does \emph{not} control
  budget, a cost-allergic CFO runs a spend freeze, legal gates on a
  data-processing agreement, and the incumbent contract auto-renews inside the
  window. Fear-selling and discounting are actively penalized; only a
  conservative, quantified case with price integrity clears the bar
  (\texttt{min\_trust}~75, \texttt{max\_discount\_pct}~5).
\end{itemize}

Together the tiers give the benchmark discrimination at both ends: an easy tier
that exposes models which cannot execute fundamentals even when handed a winnable
deal, and a hard tier that separates genuine strategic competence from
sales-shaped text.

\subsection{Design axis II: industry diversity}
\label{sec:industries}
The \numScenarios{} scenarios span \numIndustries{} distinct industries---freight
logistics, commercial banking, P\&C insurance, non-profit healthcare, omnichannel
retail, tier-1 automotive manufacturing, DTC wellness, B2B software/people
analytics, and residential home services. Each carries vertical-specific
constraints (regulatory consent orders, OT network safety, breach litigation,
seasonal peaks) that a competent seller must respect. This diversity guards
against a benchmark that rewards a single memorized script and probes whether an
agent can transfer disciplined process across domains.

\subsection{Design axis III: methodology-controlled tracks}
\label{sec:tracks}
Every scenario is evaluated in two tracks that differ in exactly one thing. In
the \textbf{out-of-box (OOB)} track the seller model receives only the deal brief
and the environment interface. In the \textbf{pack} track the identical model,
under identical seeds against the identical stochastic buyer, additionally
receives the Value Engine sales methodology as a system prompt. Because
everything except the presence of the methodology is held fixed, the paired
difference in outcomes is a near-causal estimate of \emph{methodology lift}. This
directly answers a question practitioners actually ask---does encoding a selling
methodology change agent \emph{behavior}, or only its vocabulary?---and it is, to
our knowledge, novel among agentic benchmarks. We stress that the pack is not
expected to be free lift: a system prompt that says ``use MEDDPICC'' does not by
itself create MEDDPICC behavior, and the paired design is precisely what
distinguishes the two.

\subsection{Stochastic buyers and matched pairs}
\label{sec:stochastic}
The simulated committee is instantiated by a language model whose sampling is
keyed to a per-cell seed; the seed is the label, and a cell's score is averaged
over \numSeeds{} seeds. Two properties follow. First, \textbf{robustness to
frozen-replay gaming}: because the interlocutor is re-sampled rather than replayed
from a fixed script, an agent cannot overfit a single deterministic buyer, and
reported scores reflect expected performance against a distribution of buyer
behaviors. Second, \textbf{matched win/loss pairs}: holding scenario, model, and
track fixed, two seeds can diverge only in the buyer's stochastic choices,
isolating the trajectory decisions that flipped the outcome. These pairs are the
natural unit of preference learning (Section~\ref{sec:related}) and make VEB a
data \emph{generator}, not just a scoreboard.
